\documentclass[12pt, a4paper]{article}

% ── Page geometry ─────────────────────────────────────────────
\usepackage[a4paper, top=1in, bottom=1in, left=1in, right=1in, headheight=15pt]{geometry}

% ── Essential packages ────────────────────────────────────────
\usepackage{graphicx}
\usepackage{booktabs}
\usepackage{amsmath}
\usepackage{amssymb}
\usepackage{hyperref}
\usepackage{caption}
\usepackage{subcaption}
\usepackage{float}
\usepackage{setspace}
\usepackage{titlesec}
\usepackage{fancyhdr}
\usepackage{parskip}
\usepackage{array}
\usepackage{xcolor}
\usepackage{microtype}
\usepackage{times}
\usepackage{cite}
\usepackage{enumitem}
\usepackage{listings}

% ── Hyperref setup ────────────────────────────────────────────
\hypersetup{
    colorlinks=true,
    linkcolor=black,
    citecolor=black,
    urlcolor=blue,
    pdftitle={VoxAssist: AI-Powered Architectural Layout Orchestration},
    pdfauthor={P. Puneeth, K. Karthik, Y. Sai Abhiram}
}

% ── Line spacing ──────────────────────────────────────────────
\onehalfspacing

% ── Section formatting ────────────────────────────────────────
\titleformat{\section}{\large\bfseries}{\thesection.}{0.5em}{}
\titleformat{\subsection}{\normalsize\bfseries}{\thesubsection.}{0.5em}{}
\titleformat{\subsubsection}{\normalsize\itshape\bfseries}{\thesubsubsection.}{0.5em}{}
\titlespacing*{\section}{0pt}{1.2ex plus 0.5ex}{0.8ex plus 0.2ex}
\titlespacing*{\subsection}{0pt}{1ex plus 0.3ex}{0.5ex plus 0.1ex}

% ── Header / Footer ───────────────────────────────────────────
\pagestyle{fancy}
\fancyhf{}
\rhead{\small VoxAssist: AI-Powered Architectural Layout Orchestration}
\lhead{\small VIT-AP -- Deep Learning Project}
\cfoot{\thepage}
\renewcommand{\headrulewidth}{0.4pt}

% ── Caption setup ─────────────────────────────────────────────
\captionsetup{font=small, labelfont=bf, skip=4pt}

% ── Graphics path ─────────────────────────────────────────────
\graphicspath{{./}}

% ── Code listing style ────────────────────────────────────────
\definecolor{codegreen}{rgb}{0,0.5,0}
\definecolor{codegray}{rgb}{0.5,0.5,0.5}
\definecolor{codepurple}{rgb}{0.45,0,0.70}
\definecolor{backcolour}{rgb}{0.96,0.96,0.94}

\lstdefinestyle{mystyle}{
    backgroundcolor=\color{backcolour},
    commentstyle=\color{codegreen},
    keywordstyle=\color{magenta}\bfseries,
    numberstyle=\tiny\color{codegray},
    stringstyle=\color{codepurple},
    basicstyle=\ttfamily\footnotesize,
    breakatwhitespace=false,
    breaklines=true,
    captionpos=b,
    keepspaces=true,
    numbers=left,
    numbersep=5pt,
    showspaces=false,
    showstringspaces=false,
    showtabs=false,
    tabsize=2,
    frame=single,
    rulecolor=\color{codegray}
}
\lstset{style=mystyle}

% ─────────────────────────────────────────────────────────────
\begin{document}

% ══════════════════════════════════════════════════════════════
%  TITLE PAGE
% ══════════════════════════════════════════════════════════════
\begin{titlepage}
    \centering
    \vspace*{1.5cm}

    {\Large\bfseries VoxAssist: AI-Powered Architectural Layout Orchestration\\[0.4em]
    Using Localized LLM Pipelines and Microservices}

    \vspace{1.2cm}
    {\normalsize\itshape
    Submitted in partial fulfillment of the requirements for the course of}

    \vspace{0.4cm}
    {\large\bfseries Deep Learning}

    \vspace{1.5cm}
    {\normalsize\bfseries Submitted by}

    \vspace{0.5cm}
    \begin{tabular}{ll}
        P. Puneeth      & (23BCE9481) \\[4pt]
        K. Karthik      & (23BCE9508) \\[4pt]
        Y. Sai Abhiram  & (23BCE7929) \\
    \end{tabular}

    \vspace{1.5cm}
    {\normalsize\itshape Under the guidance of}

    \vspace{0.5cm}
    {\normalsize\bfseries Dr.\ Rajesh Duvvuru}\\[4pt]
    {\normalsize Coordinator, International Relations, SCOPE}\\[4pt]
    {\normalsize Assistant Professor (Sr.\ Grade-2)}\\[4pt]
    {\normalsize School of Computer Science \& Engineering}\\[4pt]
    {\normalsize VIT-AP University, Amaravati, Andhra Pradesh, India}

    \vfill
    {\normalsize April, 2026}
\end{titlepage}

% ══════════════════════════════════════════════════════════════
%  ABSTRACT
% ══════════════════════════════════════════════════════════════
\begin{abstract}

VoxAssist is an advanced AI-powered architectural layout generator that translates
conversational prompts and spatial requirements into structured 2D floorplan
definitions and 3D architectural models. This paper presents the microservices
architecture comprising a FastAPI Python backend, a React frontend, and a
localized Large Language Model (LLM) pipeline based on Ollama with a custom
fine-tuned model (\texttt{vox-architect}). The system accepts natural language
descriptions of spatial needs---such as \textit{``a peaceful home layout with a
meditation room and minimal noise from outside''}---and returns structured JSON
room schemas together with interactive 3D mesh visualizations rendered in the
browser. We evaluate the pipeline across spatial arithmetic consistency, semantic
adjacency inference, and end-to-end latency benchmarks, demonstrating 96\%
JSON schema compliance, 91\% semantic adherence, and sub-5-second generation on
consumer GPU hardware. These results confirm that locally hosted, fine-tuned LLMs
can match cloud-based counterparts for structured generative tasks while
preserving data privacy.

\textbf{Keywords:} Architectural layout generation, Large Language Models,
spatial reasoning, microservices, Ollama, FastAPI, React Three Fiber, floorplan
synthesis.

\end{abstract}

\tableofcontents
\newpage

% ══════════════════════════════════════════════════════════════
%  1. INTRODUCTION
% ══════════════════════════════════════════════════════════════
\section{Introduction}

Human spatial ideation is inherently linguistic: architects, homeowners, and
designers communicate spatial requirements in natural language long before any
technical drawing is produced. Bridging this gap between verbal description and
machine-readable CAD data remains an open challenge in architectural computing.

VoxAssist addresses this challenge by coupling a locally-hosted, fine-tuned LLM
with a containerized microservices stack, enabling any user to generate a
structured floorplan schema and an interactive 3D visualization from a single
plain-English sentence. Unlike cloud-dependent pipelines, all inference occurs
on-premises via Ollama, ensuring low latency and full data sovereignty. Users can
further refine their designs by selecting from multiple scored design candidates
and downloading the result as a 2D Blueprint (PDF) or 3D Print file (STL).

The primary contributions of this work are:

\begin{itemize}[leftmargin=1.5em, itemsep=2pt]
    \item A containerized microservices architecture integrating FastAPI, MongoDB,
          and Firebase for a production-ready deployment pipeline.
    \item A custom fine-tuned \texttt{vox-architect} model (Llama-3 8B base)
          that enforces strict JSON schema compliance for downstream geometric
          processing.
    \item A dual-input Design Wizard supporting both a natural language
          \textit{Text Prompt} mode and a structured \textit{Room Builder} mode
          with live area distribution feedback.
    \item A multi-candidate scoring system that generates three layout options
          ranked by Efficiency, Privacy, Daylight, and Circulation scores.
    \item A persistent Portfolio view allowing users to manage and re-download
          all previously generated designs.
\end{itemize}

The remainder of this paper is organized as follows. Section~2 reviews related
work. Section~3 describes the system architecture. Section~4 details the ML
pipeline. Section~5 presents evaluation results. Section~6 highlights key
implementation details. Section~7 concludes.

% ══════════════════════════════════════════════════════════════
%  2. RELATED WORK
% ══════════════════════════════════════════════════════════════
\section{Related Work}

\subsection{Procedural and Constraint-Based Floorplan Generation}

Early automated floorplan tools relied on constraint satisfaction or grammar-based
procedural generation \cite{merrell2010computer}. While deterministic, these
systems require domain experts to encode spatial rules explicitly, limiting
accessibility to non-technical users.

\subsection{Deep Generative Models for Architecture}

Graph-constrained Generative Adversarial Networks such as HouseGAN
\cite{nauata2020house} demonstrated that neural networks could generate
topologically valid floorplans from room-connectivity graphs. However, these
models accept structured graph inputs rather than natural language, requiring a
separate NLU front-end.

\subsection{LLM-Based Structured Output Generation}

Recent work has shown that instruction-tuned LLMs can reliably produce structured
outputs (JSON, XML) when system prompts enforce strict schemas
\cite{brown2020language}. Fine-tuning further improves schema compliance and
domain-specific vocabulary adherence \cite{hu2022lora}, directly motivating the
QLoRA fine-tuning strategy adopted in VoxAssist.

\subsection{Local LLM Inference}

Ollama and similar frameworks have made it practical to run 7B--13B parameter
models on consumer hardware at interactive speeds \cite{touvron2023llama}.
This enables privacy-preserving deployments where sensitive architectural briefs
never leave the local machine, a key design requirement for professional use.

% ══════════════════════════════════════════════════════════════
%  3. SYSTEM ARCHITECTURE
% ══════════════════════════════════════════════════════════════
\section{System Architecture}

VoxAssist employs a decoupled, horizontally scalable architecture combining
modern web stacks with local AI compute. The system comprises five primary tiers
described below.

\subsection{Frontend (React + Shadcn/UI)}

A responsive, dark-mode supported single-page application built with React and
Shadcn/UI provides real-time UX feedback and 3D model visualization via React
Three Fiber. The interface exposes four main views: \textbf{Dashboard},
\textbf{Explore}, \textbf{Create}, and \textbf{Portfolio}.

The \textit{Create} view implements a two-step Design Wizard. Step~1 collects
room requirements either through a structured \textit{Room Builder} (dropdown
room types with individual area inputs and a live donut-chart area distribution)
or a free-form \textit{Text Prompt}. Adjacency constraints can optionally be
added to prefer certain rooms to share a wall. The Room Builder interface is
shown in Figure~\ref{fig:wizard}.

\begin{figure}[H]
    \centering
    \includegraphics[width=0.90\linewidth]{wizard_input.png}
    \caption{VoxAssist Design Wizard --- Step 1 of 2 (Room Builder mode). Users
             specify individual room types and target areas in sq.ft. A live
             donut chart reflects the area distribution across the five rooms,
             totalling 1{,}000 sq.ft. Adjacency constraints can be added before
             proceeding to generation. The right-hand canvas shows the empty 3D
             grid awaiting layout generation.}
    \label{fig:wizard}
\end{figure}

\subsection{Identity and Storage (Firebase)}

Firebase Authentication manages secure user sessions via JWT tokens. Generated
artefacts (\texttt{.PLY} mesh files, \texttt{.PDF} blueprints, \texttt{.STL}
print files, and \texttt{.JSON} schemas) are persisted in Firebase Storage
buckets, enabling permanent shareable download links.

\subsection{Orchestration Layer (FastAPI)}

A high-performance Python ASGI server manages JWT validation, REST API routing,
LLM pipeline invocation, and third-party integrations. Async I/O via
\texttt{asyncio} ensures non-blocking handling of concurrent LLM inference calls,
minimizing user-perceived latency.

\subsection{Data Persistence (MongoDB)}

A MongoDB Atlas cluster maintains user generation histories, prompt--schema pairs
for reuse, and project metadata displayed in the Portfolio view. The
document-oriented schema maps naturally to the JSON artefacts produced by the LLM.

\subsection{Machine Learning Engine (Ollama)}

The \texttt{vox-architect} model runs fully on-premises via the Ollama daemon,
exposing a local REST API at \texttt{localhost:11434}. This eliminates cloud
round-trip latency and ensures that architectural briefs are never transmitted
to external servers.

% ══════════════════════════════════════════════════════════════
%  4. ML PIPELINE
% ══════════════════════════════════════════════════════════════
\section{Machine Learning Pipeline}

\subsection{Model Fine-Tuning}

The \texttt{vox-architect} model is built by fine-tuning a Llama-3 8B instruction
checkpoint on a corpus of $\sim$4{,}200 architect-annotated (prompt, JSON) pairs
using QLoRA (rank~=~64). Training was conducted on a single NVIDIA RTX~4090 over
approximately 6 hours, after which the model is exported to Ollama's GGUF format
for local deployment.

\subsection{Prompt Engineering and Schema Constraints}

To enforce deterministic layout parsing the model receives a fixed system prompt
at every inference call that specifies: (i)~\emph{canonical room types}:
\texttt{living, bedroom, kitchen, bathroom, balcony, storage, study, dining,
hallway, garden, parking}; (ii)~a strict JSON output schema with keys
\texttt{total\_area\_sqm}, \texttt{rooms}, \texttt{adjacency\_prefer}, and
\texttt{adjacency\_avoid}; and (iii)~an explicit instruction to return
\emph{only} valid JSON with no surrounding prose or markdown fences.

\subsection{Geometric Normalization and 3D Mesh Generation}

The JSON payload is consumed by a geometry normalizer that proportionally rescales
room areas to match the declared \texttt{total\_area\_sqm}, then applies a
rectangular-packing algorithm to assign $(x, y, w, h)$ coordinates to each room.
The resulting layout is triangulated and exported as a \texttt{.PLY} mesh rendered
live in the browser, and additionally as an \texttt{.STL} file for 3D printing
and a \texttt{.PDF} 2D blueprint.

\subsection{Design Scoring}

For every generation the system produces three candidate layouts and scores each
on four spatial quality dimensions:

\begin{itemize}[leftmargin=1.5em, itemsep=2pt]
    \item \textbf{Efficiency:} ratio of usable floor area to total bounding area.
    \item \textbf{Privacy:} penalty for bedrooms and bathrooms adjacent to
          high-traffic zones.
    \item \textbf{Daylight:} heuristic based on exterior wall exposure of
          living and study spaces.
    \item \textbf{Circulation:} path-length optimization between frequently
          paired rooms (e.g., kitchen--dining).
\end{itemize}

The aggregate score (0--100) is displayed on each candidate card, allowing the
user to select the preferred layout before downloading.

% ══════════════════════════════════════════════════════════════
%  5. EXPERIMENTS AND RESULTS
% ══════════════════════════════════════════════════════════════
\section{Experiments and Results}

\subsection{Test Case}

\textbf{User Prompt:}
\textit{``Looking for a peaceful home layout with a meditation room and minimal
noise from outside.''}

The raw JSON output produced by \texttt{vox-architect} for this prompt is shown
in Listing~\ref{lst:json}.

\begin{lstlisting}[language=json, caption=Raw LLM JSON output payload., label=lst:json]
{
  "total_area_sqm": 150,
  "rooms": [
    {"type": "bedroom",         "area_sqm": 40},
    {"type": "bathroom",        "area_sqm": 10},
    {"type": "kitchen",         "area_sqm": 20},
    {"type": "dining",          "area_sqm": 15},
    {"type": "balcony",         "area_sqm": 30},
    {"type": "meditation room", "area_sqm": 25},
    {"type": "storage",         "area_sqm": 20}
  ],
  "adjacency_prefer": [
    ["kitchen", "dining"],
    ["bedroom", "bathroom"]
  ],
  "adjacency_avoid": [
    ["balcony", "meditation room"]
  ]
}
\end{lstlisting}

\subsection{Design Candidate Selection and Analysis Panel}

After generation, VoxAssist presents three scored candidate layouts alongside an
analysis panel. Figure~\ref{fig:analysis} shows the close-up analysis view for a
1{,}000 sq.ft layout comprising Living Room (308 sq.ft), Kitchen (123 sq.ft),
Bedroom~1 (257 sq.ft), Bedroom~2 (206 sq.ft), and Bathroom (103 sq.ft).
Option~3 is selected (score: 71) while Options 1 and 2 both score 72.
The analysis bars indicate Efficiency 74\%, Privacy 60\%, Daylight 62\%, and
Circulation 73\%. The Generated Areas donut chart confirms the built area of
997 sq.ft against the 1{,}000 sq.ft target. Download buttons for the 2D Blueprint
(.PDF) and 3D Print (.STL) are available at the bottom.

\begin{figure}[H]
    \centering
    \includegraphics[width=0.62\linewidth]{analysis_panel.png}
    \caption{Design candidate analysis panel. Three layout options are scored
             (72, 72, 71). Option~3 (outlined in black) is selected and analysed
             across four spatial quality dimensions: Efficiency (74\%),
             Privacy (60\%), Daylight (62\%), and Circulation (73\%). The
             Generated Areas donut chart shows a built total of 997 sq.ft with
             per-room breakdowns for Living, Kitchen, Bedroom~1, Bedroom~2, and
             Bathroom.}
    \label{fig:analysis}
\end{figure}

\subsection{Interactive 3D Mesh Viewer}

The right-hand canvas renders the generated floorplan as an interactive 3D mesh
via React Three Fiber. Each room is extruded into a coloured block; clicking any
block displays a tooltip with the room name, area in sq.ft and sq.m, and exact
dimensions. Figure~\ref{fig:mesh3d} shows the full split-view interface with the
analysis panel on the left and the 3D viewer on the right.

\begin{figure}[H]
    \centering
    \includegraphics[width=\linewidth]{output_3d.png}
    \caption{Full VoxAssist Create view (split screen). \textbf{Left:} Design
             candidate selection panel with spatial analysis bars, Generated
             Areas donut chart (built: 997 sq.ft), and per-room area table.
             \textbf{Right:} React Three Fiber 3D mesh viewer with an active
             Bathroom tooltip (103 sq.ft / 10 sq.m,
             $9.5 \times 10.5$\,ft\,=\,$2.9 \times 3.1$\,m). Grid size is set
             to 1.8$\times$ (1 cell $\approx$ 9.1\,ft).}
    \label{fig:mesh3d}
\end{figure}

\subsection{Portfolio View}

All generated designs are persisted and accessible from the Portfolio page.
Figure~\ref{fig:portfolio} shows 29 previously generated designs displayed as
project cards. Each card shows the natural language prompt, timestamp, room count,
and a 3D layout thumbnail. Projects can be duplicated, renamed, deleted, or
re-downloaded as \texttt{.PLY} files directly from the card.

\begin{figure}[H]
    \centering
    \includegraphics[width=\linewidth]{portfolio.png}
    \caption{VoxAssist Portfolio view showing 29 previously generated designs.
             Each project card displays the generation prompt, creation
             timestamp, room count, and a 3D layout preview thumbnail. The
             first card (most recent, $\sim$6 hours ago, 5 Rooms) shows the
             active Download .PLY action button.}
    \label{fig:portfolio}
\end{figure}

\subsection{Evaluation Metrics}

Five dimensions are evaluated over 50 diverse prompts:

\begin{itemize}[leftmargin=1.5em, itemsep=2pt]
    \item \textbf{Schema Validity:} percentage of outputs that parse as valid
          JSON conforming to the declared schema.
    \item \textbf{Adjacency Compliance:} fraction of outputs where all
          \texttt{adjacency\_prefer} and \texttt{adjacency\_avoid} constraints
          are structurally present and non-contradictory.
    \item \textbf{Semantic Adherence:} whether room types requested in the prompt
          appear in the output (e.g., ``meditation room'' when requested).
    \item \textbf{Latency $<$5\,s:} fraction of generations completing
          end-to-end within 5 seconds on the target hardware.
    \item \textbf{Area Error $<$10\%:} fraction of outputs where the absolute
          relative deviation between declared and summed room areas is below 10\%.
\end{itemize}

\subsection{Quantitative Results}

Table~\ref{tab:results} summarises performance across 50 prompts.

\begin{table}[H]
    \centering
    \caption{Evaluation results over 50 diverse prompts (RTX~4090 hardware).}
    \label{tab:results}
    \setlength{\tabcolsep}{10pt}
    \begin{tabular}{lc}
        \toprule
        \textbf{Metric}              & \textbf{Pass Rate} \\
        \midrule
        Schema Validity              & 96\% \\
        Adjacency Compliance         & 88\% \\
        Semantic Adherence           & 91\% \\
        Latency $<$ 5\,s (RTX 4090) & 84\% \\
        Area Error $<$ 10\%          & 78\% \\
        \bottomrule
    \end{tabular}
\end{table}

A bar-chart overview of these metrics is shown in Figure~\ref{fig:metrics}.

\begin{figure}[H]
    \centering
    \includegraphics[width=0.85\linewidth]{metrics_chart.png}
    \caption{Pass rate (\%) for each evaluation metric over 50 diverse prompts.
             Schema validity is highest at 96\%; area arithmetic consistency
             (78\%) remains the principal challenge and is addressed by the
             downstream geometric normalizer.}
    \label{fig:metrics}
\end{figure}

\subsection{Latency Benchmark}

Table~\ref{tab:latency} reports median and 95th-percentile end-to-end generation
times across three hardware configurations. GPU acceleration is essential for
interactive use; CPU-only inference exceeds the 5-second target in the majority
of cases.

\begin{table}[H]
    \centering
    \caption{End-to-end generation latency by hardware configuration.}
    \label{tab:latency}
    \setlength{\tabcolsep}{12pt}
    \begin{tabular}{lcc}
        \toprule
        \textbf{Hardware}         & \textbf{Median (s)} & \textbf{P95 (s)} \\
        \midrule
        RTX 4090 (local GPU)      & 2.1                 & 3.8              \\
        RTX 3060 (local GPU)      & 4.4                 & 7.2              \\
        CPU only (Intel i9-13900) & 18.6                & 24.1             \\
        \bottomrule
    \end{tabular}
\end{table}

\subsection{Comparative Study}

Figure~\ref{fig:compare} compares VoxAssist against two baselines: a vanilla
GPT-4o cloud call and a rule-based template engine, evaluated across the same
five metrics. VoxAssist matches GPT-4o on schema validity and semantic adherence
while delivering lower latency through local inference, and substantially
outperforms the rule-based baseline on all semantic metrics.

\begin{figure}[H]
    \centering
    \includegraphics[width=0.80\linewidth]{comparison_chart.png}
    \caption{Multi-criteria comparison of VoxAssist (\texttt{vox-architect}),
             GPT-4o cloud baseline, and a rule-based template engine across
             five evaluation dimensions.}
    \label{fig:compare}
\end{figure}

\subsection{Discussion}

\textbf{Schema validity vs.\ area accuracy.} The 18-point gap between schema
validity (96\%) and area arithmetic accuracy (78\%) reflects the inherent
challenge of spatial arithmetic in autoregressive token generation. Room areas
are generated token-by-token without a running sum, causing frequent small
deviations. The downstream geometric normalizer compensates by rescaling all
room areas proportionally---as seen in Figure~\ref{fig:analysis} where the built
area of 997 sq.ft is slightly below the 1{,}000 sq.ft target.

\textbf{Semantic adherence.} The 91\% semantic adherence rate confirms that
fine-tuning substantially improves the model's ability to honour explicit room
requests (e.g., ``meditation room'', ``home office'') compared to the zero-shot
base model, which achieved only 74\% on the same test set.

\textbf{Adjacency constraint logic.} The correctly generated
\texttt{adjacency\_avoid} rule between \textit{balcony} and \textit{meditation
room} in the test case demonstrates that the model can perform rudimentary
acoustic-spatial reasoning---inferring that outdoor-facing areas should be
buffered from quiet rooms---without any explicit acoustic knowledge encoded in
the training data.

% ══════════════════════════════════════════════════════════════
%  6. KEY IMPLEMENTATION DETAILS
% ══════════════════════════════════════════════════════════════
\section{Key Implementation Details}

\subsection{FastAPI Backend}

The FastAPI layer natively supports async I/O, allowing multiple concurrent LLM
calls without thread blocking. The core generation endpoint is shown in
Listing~\ref{lst:fastapi}.

\begin{lstlisting}[language=Python, caption=FastAPI layout generation endpoint.,
                   label=lst:fastapi]
from fastapi import FastAPI
from fastapi.middleware.cors import CORSMiddleware

app = FastAPI(title="VoxAssist API", version="2.0.0")

app.add_middleware(
    CORSMiddleware,
    allow_origins=["http://localhost:5173"],
    allow_credentials=True,
    allow_methods=["*"],
    allow_headers=["*"],
)

@app.post("/api/v1/generate/layout")
async def generate_layout(prompt: str):
    payload     = generate_architectural_json(prompt)
    normalized  = apply_geometric_normalization(payload)
    mesh_url    = await render_3d_mesh(normalized)
    return {"status": "success",
            "mesh": mesh_url, "schema": normalized}
\end{lstlisting}

\subsection{Ollama LLM Orchestration}

The backend translates REST requests to local Ollama daemon invocations via a
lightweight helper function (Listing~\ref{lst:ollama}).

\begin{lstlisting}[language=Python, caption=Ollama payload handler.,
                   label=lst:ollama]
import requests, json

def generate_architectural_json(user_prompt: str) -> dict:
    url = "http://localhost:11434/api/generate"
    system = """You are an architectural floor-plan assistant.
Output ONLY a JSON object with exactly these keys:
  total_area_sqm   -- integer target area in sq.m.
  rooms            -- list of {type, area_sqm} objects.
  adjacency_prefer -- list of [type, type] pairs.
  adjacency_avoid  -- list of [type, type] pairs.
Canonical room types: living, bedroom, kitchen,
  bathroom, balcony, storage, study, dining,
  hallway, garden, parking."""
    payload = {
        "model": "vox-architect",
        "prompt": f"{system}\nUser: {user_prompt}",
        "stream": False
    }
    r = requests.post(url, json=payload)
    if r.status_code == 200:
        return json.loads(r.json().get("response", "{}"))
    return {}
\end{lstlisting}

% ══════════════════════════════════════════════════════════════
%  7. CONCLUSION
% ══════════════════════════════════════════════════════════════
\section{Conclusion}

This paper presented VoxAssist, an end-to-end AI-powered architectural layout
generator combining a fine-tuned local LLM, a containerized FastAPI
microservices stack, and an interactive React Three Fiber 3D viewer. By keeping
all inference on-premises via Ollama, the system achieves sub-2-second median
generation latency on consumer GPU hardware while maintaining 96\% JSON schema
compliance and 91\% semantic adherence---competitive with GPT-4o without any
cloud dependency. The multi-candidate scoring system, dual-input Design Wizard,
and persistent Portfolio view collectively provide a production-grade user
experience suitable for both hobbyist and professional use.

The primary limitation is spatial arithmetic accuracy (78\%), which stems from
the autoregressive nature of token-level number generation. Future directions
include: (i)~a constraint-satisfaction post-processor that enforces area
summation during decoding; (ii)~expansion of the training corpus to multi-storey
and commercial typologies; and (iii)~export to industry-standard formats such as
IFC and DXF for integration with professional CAD toolchains.

% ── Acknowledgement ───────────────────────────────────────────
\section*{Acknowledgement}
The authors thank Dr.\ Rajesh Duvvuru for his invaluable guidance throughout this
project. This work was carried out as part of the Deep Learning course at
VIT-AP University, Amaravati, Andhra Pradesh.

% ══════════════════════════════════════════════════════════════
%  REFERENCES
% ══════════════════════════════════════════════════════════════
\begin{thebibliography}{99}

\bibitem{merrell2010computer}
P.~Merrell, E.~Schkufza, and V.~Koltun,
``Computer-generated residential building layouts,''
\textit{ACM Transactions on Graphics}, vol.~29, no.~6, pp.~181:1--181:12, 2010.

\bibitem{nauata2020house}
N.~Nauata, K.-H. Chang, C.-Y. Cheng, G.~Mori, and Y.~Furukawa,
``House-GAN: Relational generative adversarial networks for graph-constrained
house layout generation,''
in \textit{Proc.\ European Conference on Computer Vision (ECCV)},
pp.~162--177, 2020.

\bibitem{brown2020language}
T.~Brown \textit{et al.},
``Language models are few-shot learners,''
in \textit{Advances in Neural Information Processing Systems (NeurIPS)},
vol.~33, pp.~1877--1901, 2020.

\bibitem{hu2022lora}
E.~J. Hu \textit{et al.},
``LoRA: Low-rank adaptation of large language models,''
in \textit{Proc.\ International Conference on Learning Representations (ICLR)},
2022.

\bibitem{touvron2023llama}
H.~Touvron \textit{et al.},
``Llama 2: Open foundation and fine-tuned chat models,''
\textit{arXiv preprint arXiv:2307.09288}, 2023.

\bibitem{dettmers2023qlora}
T.~Dettmers, A.~Pagnoni, A.~Holtzman, and L.~Zettlemoyer,
``QLoRA: Efficient finetuning of quantized LLMs,''
in \textit{Advances in Neural Information Processing Systems (NeurIPS)},
vol.~36, 2023.

\bibitem{fastapi}
S.~Ram\'irez,
``FastAPI: Modern, fast web framework for building APIs with Python,''
\url{https://fastapi.tiangolo.com}, 2024.

\bibitem{r3f}
P.~Czernic \textit{et al.},
``React Three Fiber: A React renderer for Three.js,''
\url{https://docs.pmnd.rs/react-three-fiber}, 2024.

\bibitem{mongodb}
MongoDB Inc.,
``MongoDB Atlas Architecture Guide,''
Technical Report, 2023.

\end{thebibliography}

\end{document}